\section{Analytic Functions and Taylor Series}
\index{function!analytic}
%%%%%%%%%%%%%%%%%%%%%%%%%%%%%%%%%%%%%%
%%%%%%%%%%%%%%%%%%%%%%%%%%%%%%%%%%%%%%
%%%%%%%%%%%%%%%%%%%%%%%%%%%%%%%%%%%%%%

In the chapter on power series, we studied power series of the form
\begin{equation}\label{eq:4748219}
    g(z) = \sum_{k=0}^{+\infty} a_k z^k.
\end{equation}
It is clear that if we fix a point $z_0$, we can perform a change of variable
and observe that the series
\begin{equation}\label{eq:33783}
  f(z) = \sum_{k=0}^{+\infty} a_k \cdot(z-z_0)^k = g(z-z_0)
\end{equation}
converges in the circle $B_R(z_0) = \ENCLOSE{z \in \ZZ\colon \abs{z-z_0}< R}$
where $R$ is the radius of the power series~\eqref{eq:4748219}. We will say that
the series in equation~\eqref{eq:33783} is a power series centered at $z_0$ and
will call $R$ its radius of convergence.

It is then natural to ask whether a power series centered at a point $z_0$ can
be written as a power series centered at the point $z_1$, at least when
$\abs{z_1-z_0}<R$. The affirmative answer is given by the following.

\begin{theorem}[Translation of a Power Series]
\index{power series!translation}%
\index{translation!of a power series}%
\label{th:488456367}
Consider the power series~\eqref{eq:33783} centered at $z_0$ with radius of
convergence $R\in (0,+\infty]$, take a point $z_1 \in B_R(z_0)$ and set
$r=R-\abs{z_1-z_0}$ so that $B_r(z_1)\subset B_R(z_0)$. Then for every $z\in
B_r(z_1)$, we have
\begin{equation}\label{eq:48478643}
  f(z) = \sum_{j=0}^{+\infty} b_j (z - z_1)^j
\end{equation}
for suitable coefficients $b_j$. In particular, the power series
in~\eqref{eq:48478643}, centered at $z_1$, has a radius of convergence not less
than $r$.
\end{theorem}
%
\begin{proof}
To simplify the notation, we can assume, without loss of generality, that
$z_0=0$. If we set $h=z-z_1$, we have
\begin{align*}
f(z) 
&= \sum_{k=0}^{+\infty} a_k z^k 
= \sum_{k=0}^{+\infty} a_k (z_1+h)^k
= \sum_{k=0}^{+\infty} \sum_{j=0}^k a_k {k \choose j}z_1^{k-j}h^j.
\end{align*}
We would now like to apply theorem~\ref{th:somma_Cauchy} to exchange the order
of the two sums. To do this, we need to verify the convergence of the following
series:
\begin{align*}
  \sum_{k=0}^{+\infty} \sum_{j=0}^k \abs{a_k} {k \choose j}\abs{z_1}^{k-j}\abs{h}^j 
  &= \sum_{k=0}^{+\infty} \abs{a_k \enclose{\abs{z_1}+\abs{h}}^k}.
\end{align*}
This series converges because, by hypothesis, $\abs{z_1}+\abs{h} \le \abs{z_1} +
r < R$, and thus we are within the radius of convergence where the series
converges absolutely. Therefore, we can exchange the sums and obtain:
\[
    f(z) = \sum_{j=0}^{+\infty} \sum_{k=j}^{+\infty} a_k {k\choose j} z_1^{k-j}
  h^j
    = \sum_{j=0}^{+\infty} b_j h^j
  \qquad\text{with }
b_j = \sum_{k=j}^{+\infty}
 a_k {k\choose j}z_1^{k-j}.
\]
Thus, the power series $\sum b_j (z-z_1)^j$ has a radius of convergence not less
than $r$ and has sum $f(z)$, as we wanted to prove.
\end{proof}

The theory of analytic functions could be developed, without substantially
changing anything, for functions of a complex variable. However, it would be
necessary to introduce the concept of derivative in the complex sense, and
therefore, for simplicity, we limit ourselves to real functions, which are our
subject of study.

\begin{definition}[Analytic Function]
\mymargin{analytic function}%
\index{function!analytic}
Let $A\subset \RR$ and let $f\colon A\subset \RR \to \RR$ be a function. We will
say that $f$ is \emph{analytic} if for every $x_0\in A$, there exists $R>0$ such
that $(x_0-R,x_0+R)\subset A$ and there exists a sequence of real numbers $a_k$
such that for every $x\in (x_0-R,x_0+R)$, we have:
\[
  f(x) = \sum_{k=0}^{+\infty} a_k (x-x_0)^k.
\]
\end{definition}

Theorem~\ref{th:488456367} tells us, in particular, that the sum of a power
series is an analytic function, at least within the radius of convergence. For
example, the function $f(x)=e^x$ is an analytic function defined on all of $\RR$
since, thanks to theorem~\ref{th:serie_esponenziale}, $f(x)$ turns out to be the
sum of a power series with radius of convergence $R=+\infty$ centered at $x=0$.

\begin{theorem}[Developability in Taylor Series]
\label{th:48765}%
If $f\colon A\subset \RR\to\RR$ is an analytic function, then $f\in C^\infty(A)$
and for every $x_0\in A$, there exists $R>0$, such that for every $x \in (x_0-R,
x_0+R)\subset A$, we have
\begin{equation}\label{eq:2344494}
  f(x) = \sum_{k=0}^{+\infty} \frac{f^{(k)}(x_0)}{k!}(x-x_0)^k.
\end{equation}
Also, the derivative of $f$ is an analytic function, and the Taylor series of
the derivative is obtained by differentiating term by term the Taylor series of
$f$.
\end{theorem}
%
\begin{proof}
Without loss of generality, let us assume that $x_0=0$. Since $f$ is an analytic
function, by definition, we know that there exists $R>0$ such that for
$\abs{x}<R$, we have
\[
f(x) = \sum_{k=0}^{+\infty} a_k x^k
\]
with suitable coefficients $a_k$. We want to show that $f(x)$ is differentiable.
Informally, we would like to perform the following steps:
\begin{align*}
\frac{f(x+h) - f(x)}{h}
&= \sum_{k=1}^{+\infty} a_k \cdot \frac{(x+h)^k-x^k}{h}
= \sum_{k=1}^{+\infty} a_k \cdot \frac{\sum_{j=1}^k {k\choose j} \cdot h^j \cdot
x^{k-j}}{h} \\
&= \sum_{k=1}^{+\infty} a_k \cdot \sum_{j=1}^k {k\choose j} \cdot h^{j-1} \cdot
x^{k-j} \\
&= \sum_{k=1}^{+\infty} a_k \cdot \sum_{j=0}^{k-1} {k\choose j+1} \cdot h^j
\cdot x^{k-j-1} \\
&= \sum_{j=0}^{+\infty} \enclose{\sum_{k=j+1}^{+\infty} a_k \cdot {k\choose j+1}
\cdot x^{k-j-1}} h^j.
\end{align*}
If we set
\[
  c_{k,j} = a_k\cdot {k\choose j+1}\cdot x^{k-j-1}\cdot h^j
\]
the steps are justified if the series $\sum_{k,j} c_{k,j}$ is absolutely
convergent. But, by appropriately placing absolute values in the steps already
taken, we obtain
\[
  \sum_{j=0}^{+\infty} \sum_{k=j+1}^{+\infty} \abs{c_{k,j}}
    =  \sum_{k=1}^{+\infty} \abs{a_k} \frac{\enclose{\abs x + \abs
  h}^k-\abs{x}^k}{\abs{h}}
\]
and thus we can assert that there is absolute convergence if $\abs{x}<R$ and
$\abs{x}+\abs{h}<R$, that is, if $\abs{h}< R-\abs{x}$. In such a case, the steps
taken before are justified, and therefore the result is a power series
convergent for $\abs{h}<R-\abs{x}$. In particular, for fixed $x$, the sum of the
power series obtained at the end is a continuous function, and as $h\to 0$, it
tends to the constant term of the series (the one with $j=0$), thus we have
\[
  \lim_{h\to 0} \frac{f(x+h)-f(x)}{h}
  = \sum_{k=1}^{+\infty} a_k \cdot {k \choose 1} \cdot x^{k-1}
  = \sum_{k=1}^{+\infty} k a_k \cdot x^{k-1}.
\]
We have found that the derivative of the power series is equal to the power
series of the derivatives.

Since the series of derivatives has the same radius of convergence $R$ as the
original series (theorem~\ref{th:raggio_serie_derivate}), the procedure can be
iterated, finding that each derivative is differentiable, and for every $x\in
(-R,R)$, we will have:
\[
    f^{(n)}(x) = \sum_{k=n}^{+\infty}k (k-1) \cdots (k-n+1) \cdot a_k \cdot
  x^{k-n}.
\]
In particular,
\[
  f^{(n)}(0) = n! \cdot a_n
\]
and therefore
\[
  f(x) = \sum_{k=0}^{+\infty} \frac{f^{(k)}(0)}{k!} x^k
\]
as we needed to prove.
\end{proof}

We have seen that analytic functions are of class $C^\infty$. The following
example shows us that the converse is not true, and thus the class of analytic
functions is strictly contained within the class of $C^\infty$ functions.

\begin{example}[A $C^\infty$ Function That Is Not Analytic]
The function
\[
  f(x) =
  \begin{cases}
    e^{-\frac 1 {x^2}} & \text{if $x\neq 0$}\\
    0 & \text{if $x=0$}
  \end{cases}
\]
is of class $C^\infty$ because for $x\neq 0$, we can calculate all the
derivatives and observe that they can be written in the form
\[
  f^{(n)}(x) = \frac{P_n(x)}{x^{3n}}e^{-\frac 1 {x^2}}
\]
for a suitable polynomial $P_n$ (this can be proven by induction). Thus, as
$x\to 0$, we have, for every $n\in \NN$,
\[
  f^{(n)}(x) \to 0
\]
and therefore the function $f$ is infinitely differentiable even at the point
$x=0$ (thanks to proposition~\ref{prop:4384774}), and it results in $f^{(n)}(0)
= 0$ for every $n\in \NN$. If the function $f$ were analytic in a neighborhood
of $0$, we would have, by theorem~\ref{th:48765}:
\[
  f(x) = \sum_{k=0}^{+\infty} 0\cdot x^k = 0
\]
which is absurd since $f(x)>0$ for every $x\neq 0$.
\end{example}

\begin{theorem}[Binomial Series]
\label{th:serie_binomiale}%
\mymargin{binomial series}%
\index{series!binomial}\index{binomial!series}
For every $\alpha\in \RR$ and for every $x\in (-1,1)$, we have
\[
  (1+x)^{\alpha}  = \sum_{k=0}^{+\infty} {\alpha \choose k} \cdot x^k.
\]
\end{theorem}
%
\begin{proof}
First, observe that the power series with coefficient $a_k = {\alpha \choose k}$
has a radius of convergence $R=1$. Indeed, using the ratio test, we find:
\[
  \frac{\abs{a_{n+1}}}{\abs{a_n}}
  = \frac{\abs{\alpha (\alpha-1) \cdots (\alpha-n)}}{(n+1)!}
  \cdot \frac{n!}{\abs{\alpha(\alpha-1) \cdots (\alpha-n+1)}}
  = \frac{\abs{\alpha-n}}{n+1} \to 1.
\]
Thus, the function
\[
 f(x) = \sum_{k=0}^{+\infty} {\alpha \choose k} x^k
\]
is defined for every $x\in(-1,1)$. Thanks to theorem~\ref{th:488456367}, the
function $f$ is developable in a power series in the neighborhood of every point
of the interval $(-1,1)$ and thus is an analytic function on this interval.
Moreover, its derivative has as its development the series of derivatives:
\[
  f'(x) = \sum_{k=1}^{+\infty} k {\alpha \choose k} \cdot x^{k-1}.
\]

To prove that $f(x) = (1+x)^\alpha$, it suffices to show that the ratio:
\[
  u(x) = \frac{f(x)}{(1+x)^\alpha}
\]
is constantly equal to $1$. By direct verification, we see that $f(0)=1$, so we
have at least $u(0) = 1$. It is then enough to show that $u$ is constant, i.e.,
that $u'=0$. But we have
\begin{align*}
  u'(x)
    &= \frac{f'(x) (1+x)^\alpha - f(x) \alpha (1+x)^{\alpha-1}}{(1+x)^{2\alpha}}
  \\
  &= \frac{(1+x) f'(x) - \alpha f(x)}{(1+x)^{\alpha+1}}.
\end{align*}
We then observe that
\begin{align*}
(1+x)f'(x)
&= \sum_{k=1}^{+\infty} k{\alpha \choose k} x^{k-1}
+ \sum_{k=1}^{+\infty} k{\alpha \choose k} x^{k}\\
&= \sum_{k=0}^{+\infty} (k+1){\alpha \choose k+1} x^k
+ \sum_{k=0}^{+\infty} k {\alpha \choose k} x^k \\
&= \sum_{k=0}^{+\infty} \Enclose{\frac{\alpha(\alpha-1) \cdots (\alpha-k)}{k!}
 + k \frac{\alpha(\alpha-1)\cdots(\alpha-k+1)}{k!}} x^k \\
  &= \sum_{k=0}^{+\infty}
 \frac{\alpha(\alpha-1)\cdots(\alpha-k+1)}{k!}\Enclose{(\alpha-k)+k} x^k\\
 &= \alpha \sum_{k=0}^{+\infty} {\alpha \choose k} x^k = \alpha \cdot f(x).
\end{align*}
Thus, $u'(x)=0$ and $f(x) = (1+x)^\alpha$ for every $x\in(-1,1)$, as we wanted
to prove.
\end{proof}

It is not enough for a function to be of class $C^\infty$ to guarantee that it
is also analytic, and therefore the following can be useful.

\begin{theorem}[Criterion of Analyticity]
\label{th:criterio_analitica}
Let $f\in C^\infty(A)$ be a real function defined on an open set $A\subset \RR$.
Suppose that for every $x_0\in A$, there exist $r>0$, $M>0$, and $L>0$ such that
$[x_0-r, x_0+r] \subset A$ and for every $x\in [x_0-r, x_0+r]$ and for every
$n\in \NN$, we have
\[
  \frac{\abs{f^{(n)}(x)}}{n!} \le M\cdot L^n.
\]
Then $f$ is analytic in $A$.
\end{theorem}
%
\begin{proof}
Fix $x_0\in A$. We need to show that in a neighborhood of $x_0$, it is possible
to write $f$ as the sum of a power series $\sum a_k (x-x_0)^k$. Necessarily, by
theorem~\ref{th:48765}, the power series must be the Taylor series, i.e., we
want to show that there exists $\rho>0$ such that if $\abs{x-x_0}< \rho$, we
have
\[
  f(x) = \sum_{k=0}^{+\infty} \frac{f^{(k)}(x_0)}{k!}(x-x_0)^k.
\]
Thanks to the Taylor formula with Lagrange remainder
(theorem~\ref{th:taylor_lagrange}), we know that
\[
  f(x) = \sum_{k=0}^{n} \frac{f^{(k)}(x_0)}{k!}(x-x_0)^k
   + \frac{f^{(n+1)}(y)}{(n+1)!}(x-x_0)^{n+1}.
\]
Thus, it will be sufficient to show that for fixed $x$, as $n\to +\infty$, the
remainder tends to zero. But indeed, if $\abs{x-x_0}\le \rho$, we have
\begin{align*}
 \abs{\frac{f^{(n+1)}(y)}{(n+1)!}(x-x_0)^{n+1}}
 & \le  M\cdot L^{n+1}\abs{x-x_0}^{n+1}
 = M (L\rho)^{n+1}\to 0
\end{align*}
if we choose $\rho$ such that $\rho < 1/L$.
\end{proof}

\begin{theorem}[Analyticity of Some Elementary Functions]
  \label{th:serie_taylor}
    The functions $e^x$, $\sin x$, $\cos x$, $\ln x$ are analytic functions. The
  functions $\arcsin$ and $\arccos$ are analytic on the open interval $(-1,1)$.
  For every $\alpha \in \RR$, the function $x^\alpha$ is analytic on the
  interval $(0,+\infty)$. Moreover, the formulas reported in
  table~\ref{tb:taylor2} hold.
\end{theorem}
%
\begin{table}
\begin{align*}
e^x &= \sum_{k=0}^{+\infty} \frac{x^k}{k!}\\
\sin x &= \sum_{k=0}^{+\infty} (-1)^k\frac{x^{2k+1}}{(2k+1)!}\\
\cos x &= \sum_{k=0}^{+\infty} (-1)^k\frac{x^{2k}}{(2k)!}\\
\sinh x &= \sum_{k=0}^{+\infty} \frac{x^{2k+1}}{(2k+1)!}\\
\cosh x &= \sum_{k=0}^{+\infty} \frac{x^{2k}}{(2k)!} \\
(1+x)^\alpha &= \sum_{k=0}^{+\infty} {\alpha \choose k} x^k, \qquad \text{if
$\abs{x}<1$} \\
\ln (1+x) &= \sum_{k=1}^{+\infty} (-1)^{k-1} \frac{x^k}{k}, \qquad \text{if $x\in(-1,1]$} \\
\arctg x &= \sum_{k=0}^{+\infty} (-1)^k \frac{x^{2k+1}}{2k+1}, \qquad \text{if $x\in [-1,1]$} \\
\arcsin x &= \sum_{k=0}^{+\infty} \frac{(2k-1)!!}{(2k)!!} \frac{x^{2k+1}}{2k+1}, 
  \qquad \text{if $x\in[-1,1]$} \\
\arccos x &= \frac \pi 2 - \arcsin x
\end{align*}
\caption{Taylor series expansions of some elementary functions.
\index{Taylor!series expansions of elementary functions}%
\index{expansion!Taylor series}%
See theorem~\ref{th:serie_taylor}.}
\label{tb:taylor2}%
\end{table}
%
\begin{proof}
    That $e^x$ is analytic has already been observed (it is a consequence of
  theorems~\ref{th:serie_esponenziale} and \ref{th:488456367}). The functions
  $\sin x$ and $\cos x$ have all derivatives bounded (between $-1$ and $1$), and
  therefore, thanks to theorem~\ref{th:criterio_analitica}, they turn out to be
  analytic over all $\RR$.
  
    Regarding the function $f(x)=x^\alpha$, chosen any point $x_0>0$, we can use
  theorem~\ref{th:serie_binomiale}:
    \begin{align*}
    x^\alpha
    &= (x_0 + x-x_0)^\alpha
    = x_0^\alpha \enclose{1+\frac {x-x_0} {x_0}}^\alpha \\
        &= x_0^\alpha \sum_{k=0}^{+\infty} {\alpha \choose k}
    \enclose{\frac{x-x_0}{x_0}}^k
    = \sum_{k=0}^{+\infty} x_0^{\alpha-k}{\alpha \choose k} (x-x_0)^k
  \end{align*}
    and thus $x^\alpha$ is also developable in a power series in the
  neighborhood of any point $x_0>0$ and is therefore an analytic function.

    We now observe that if $f$ is differentiable and $f'$ is analytic, then also
  $f$ is analytic. Indeed, if
  \[
    f'(x) = \sum_{k=0}^{+\infty} a_k (x-x_0)^k
  \]
  on a certain interval $I$, then the function
  \begin{equation}\label{eq:4569903}
    g(x) = f(x_0) + \sum_{k=0}^{+\infty} \frac{a_k}{k+1}x^{k+1}
  \end{equation}
    is also analytic on $I$ (since it is defined by a power series), but clearly
  $g(x_0) = f(x_0)$ and $g'(x) = f'(x)$ for every $x\in I$, thus $g-f$ is
  constantly equal to $0$, i.e., $f(x) = g(x)$.

    The previous observation applies to the functions $\ln x$, $\arctg x$,
  $\arcsin x$, and $\arccos x$ since their derivative is a power function and
  therefore is analytic as already seen:
  \begin{align*}
    D \ln (1+x) &= \enclose{1+x}^{-1}, &
    D \arctg x &= \enclose{1+x^2}^{-1}, \\
    D \arcsin x &= \enclose{1+x^2}^{-\frac 1 2}, &
    D \arccos x &= -\enclose{1+x^2}^{-\frac 1 2}.
  \end{align*}

    We then find the series expansions reported in table~\ref{tb:taylor2}, valid
  within the radius of convergence of these series, i.e., for $x\in (-1,1)$.

    We recall, however, theorem~\ref{th:lemma_abel} (Abel's lemma), which
  guarantees that if we have a power series whose radius of convergence is $R$
  and if the series is convergent at the point $x$ with $x=\pm R$, then the sum
  of the series turns out to be a continuous function up to the point $x$. For
  example, the sum of the Taylor series of the function $\ln(1+x)$ is continuous
  at the point $x=1$ since at that point it is convergent (it is indeed the
  alternating harmonic series, convergent by the Leibniz criterion). Thus, we
  have
  \begin{equation}\label{eq:somma_serie_leibniz}
    \sum_{k=1}^{+\infty} \frac{(-1)^{k-1}} k 
    = \lim_{x\to 1^-} \sum_{k=1}^{+\infty} (-1)^{k-1} \frac{x^k}{k}
    = \lim_{x\to 1^-} \ln (1+x) = \ln 2 
  \end{equation}
    and the series expansion of the logarithm $\ln(1+x)$ is valid over the
  entire semi-open interval $(-1,1]$.

    The same holds for the functions $\arctg x$, $\arcsin x$, and $\arccos x$,
  whose series are convergent over the entire closed interval $[-1,1]$.

    Naturally, the coefficients of the Taylor series we have found now coincide
  with the coefficients of the Taylor polynomials we had already found in
  table~\ref{tb:taylor} on page \pageref{tb:taylor} since truncating the Taylor
  series yields a polynomial that satisfies the Taylor formula~\eqref{eq:Taylor}
  and is therefore precisely the Taylor polynomial.

    The Taylor series of the functions $\sinh$ and $\cosh$ is immediately
  obtained by summing or subtracting the Taylor series of $e^x$ and $e^{-x}$.
\end{proof}

The previous theorem allows us to write remarkable identities. Setting $x=1$ in
the series that defines the logarithm, we have already observed that we obtain
the sum of the alternating harmonic series of \emph{Newton-Mercator}%
\mymargin{Newton-Mercator}%
\index{Newton-Mercator}%
\index{formula!of Newton-Mercator}%
\index{series!of Newton-Mercator}%
\index{series!harmonic!alternating}%
\index{$\sqrt 2$!series expansion}%
\begin{equation}\label{eq:serie_ln2}
  \ln 2 = 1 - \frac 1 2 + \frac 1 3 - \frac 1 4 + \frac 1 5 - \frac 1 6 + \dots
\end{equation}
Setting $x=1$ in the Taylor series of the arctangent, we obtain the formula of
\emph{Gregory-Leibniz}%
\mymargin{Gregory-Leibniz}%
\index{Gregory-Leibniz}
\index{$\pi$!Gregory-Leibniz formula}%
\index{Gregory!approximation $\pi$}%
\index{Leibniz!approximation $\pi$}%
\index{formula!of Gregory-Leibniz for $\pi/4$}%
\[
  \frac \pi 4 =
   1 - \frac 1 3 + \frac 1 5 - \frac 1 7 + \frac 1 9 - \frac 1 {11} + \dots
\]

Even the series expansion of the arcsine determines a useful formula for
calculating the decimal digits of $\pi$. For faster convergence, it is advisable
to take $x$ as close to zero as possible; in the following exercise, we take
$x=\frac 1 2$ to find the well-known approximation $\pi \approx 3.14$.

\begin{exercise}[Calculating Digits of $\pi$]
  \label{ex:cifre_pi}
%  We can refine the calculation of the digits of $\pi$
%  done in exercise~\ref{ex:4899264}.
  Thanks to the Taylor expansion of the $\arcsin$ function, we know that
  \[
    \frac \pi 6
    = \arcsin \frac 1 2
        = \sum_{k=0}^{+\infty} \frac{(2k-1)!!}{(2k)!!} \frac{1}{(2k+1)\cdot
    2^{2k+1}}
  \]
  thus
  \[
    \pi = 3 \sum_{k=0}^n \frac{(2k-1)!!}{(2k)!!} \frac{1}{(2k+1)\cdot 4^k} + \eps_{n+1}
  \]
  where
  \begin{align*}
   \eps_n
   &= 3 \sum_{k=n}^{+\infty} \frac{(2k-1)!!}{(2k)!!} \frac{1}{(2k+1)\cdot 4^k}
   \le 3 \sum_{k=n}^{+\infty} \frac{1}{(2k+1)\cdot 4^k} \\
   &\le \frac{3}{(2n+1)\cdot 4^n} \sum_{k=0}^{+\infty} \frac{1}{4^k}
   = \frac{3}{(2n+1)\cdot 4^n} \cdot \frac{1}{1-\frac 1 4}
   = \frac{4}{(2n+1)\cdot 4^n}.
 \end{align*}
 For $n=3$, we thus obtain
 \begin{align*}
  \pi
    &= 3 + \frac{3}{2\cdot 3 \cdot 4} + \frac{3\cdot 3}{4\cdot 2 \cdot 5 \cdot
  4^2}
  + \frac{3\cdot 5\cdot 3}{6\cdot 4\cdot 2 \cdot 7 \cdot 4^3} + \eps_4 \\
  &= 3+\frac{1}{8} + \frac{9}{640} + \frac{45}{21504} + \eps_4
\end{align*}
 with
 \[
 0 < \eps_4 < \frac{1}{9\cdot 4^3} < 0.0018
 \]
 from which
 \[
   3.141 < \pi < 3.143
 \]

  Using this method with a calculator, we can quickly find many exact digits, as
 reported in table~\ref{tab:cifre_pi}.
\end{exercise}
%
\begin{table}
\begin{center}
\begin{tabular}{r}
\ttfamily\footnotesize 3.1415926535 8979323846 2643383279 5028841971 6939937510 \\
\ttfamily\footnotesize   5820974944 5923078164 0628620899 8628034825 3421170679 \\
\ttfamily\footnotesize   8214808651 3282306647 0938446095 5058223172 5359408128 \\
\ttfamily\footnotesize   4811174502 8410270193 8521105559 6446229489 5493038196 \\
\ttfamily\footnotesize   4428810975 6659334461 2847564823 3786783165 2712019091 \\
\ttfamily\footnotesize   4564856692 3460348610 4543266482 1339360726 0249141273 \\
\ttfamily\footnotesize   7245870066 0631558817 4881520920 9628292540 9171536436 \\
\ttfamily\footnotesize   7892590360 0113305305 4882046652 1384146951 9415116094 \\
\ttfamily\footnotesize   3305727036 5759591953 0921861173 8193261179 3105118548 \\
\ttfamily\footnotesize   0744623799 6274956735 1885752724 8912279381 8301194912 \\
\ttfamily\footnotesize   9833673362 4406566430 8602139494 6395224737 1907021798 \\
\ttfamily\footnotesize   6094370277 0539217176 2931767523 8467481846 7669405132 \\
\ttfamily\footnotesize   0005681271 4526356082 7785771342 7577896091 7363717872 \\
\ttfamily\footnotesize   1468440901 2249534301 4654958537 1050792279 6892589235 \\
\ttfamily\footnotesize   4201995611 2129021960 8640344181 5981362977 4771309960 \\
\ttfamily\footnotesize   5187072113 4999999837 2978049951 0597317328 1609631859 \\
\ttfamily\footnotesize   5024459455 3469083026 4252230825 3344685035 2619311881 \\
\ttfamily\footnotesize   7101000313 7838752886 5875332083 8142061717 7669147303 \\
\ttfamily\footnotesize   5982534904 2875546873 1159562863 8823537875 9375195778 \\
\ttfamily\footnotesize   1857780532 1712268066 1300192787 6611195909 2164201989
\end{tabular}
\end{center}
\caption{The first 1000 decimal digits of the number $\pi$ calculated using the method employed in exercise~\ref{ex:cifre_pi}. See the code on page~\pageref{code:compute_pi}.}
\label{tab:cifre_pi}
\index{$\pi$!decimal digits}
\index{digits!$\pi$}
\end{table}

\begin{theorem}[Principle of Identity of Analytic Functions]
    If $f$ and $g$ are two analytic functions defined on the same open interval
  $I\subset \RR$, then the following are equivalent:
  \begin{enumerate}
   \item $f(x)=g(x)$ for every $x\in I$,
   \item there exists $x_0\in I$ such that for every $k\in \NN$, we have $f^{(k)}(x_0) = g^{(k)}(x_0)$,
   \item there exists $x_0\in I$ and $\rho>0$ such that for every $x\in I$ with $\abs{x-x_0}<\rho$, we have $f(x)=g(x)$.
  \end{enumerate}
  \end{theorem}
  \begin{proof}
    If $f=g$, then clearly the derivatives of $f$ and $g$ coincide at all
  points, so $1\implies 2$.
  
    If $f$ and $g$ are analytic and $x_0\in I$, there exists a neighborhood of
  the point $x_0$ in which both functions can be written as the sum of their
  Taylor series. But if $f$ and $g$ have the same derivatives at the point
  $x_0$, their Taylor series coincide, and thus the functions coincide within
  the radius of convergence. Therefore, $2\implies 3$.
  
    We are left to prove that $3 \implies 1$. By hypothesis, we know that there
  exists $x_0\in I$ such that $f(x) = g(x)$ in a neighborhood of $x_0$. Suppose,
  for contradiction, that there exists at least one point $x\in I$ where $f(x)
  \neq g(x)$. Without loss of generality, we can assume that $x>x_0$ and we can
  take the infimum of such points:
  \[
    x_1 = \inf\ENCLOSE{x\in I\colon x>x_0, f(x)\neq g(x)}.
  \]
    We know that $x_1>x_0$ because in a neighborhood of $x_0$, the two functions
  coincide by hypothesis. By definition of $x_1$, we also know that $f(x)=g(x)$
  for every $x\in [x_0,x_1)$, and therefore, by differentiating, we can assert
  that $f^{(k)}(x) = g^{(k)}(x)$ for every $k\in \NN$ and for every $x\in
  [x_0,x_1)$. Since $f$ and $g$ are of class $C^\infty$, we can conclude, by
  continuity, that also $f^{(k)}(x_1)=g^{(k)}(x_1)$ for every $k\in \NN$. Thus,
  the two functions have the same derivatives at the point $x_1$. Since $f$ and
  $g$ are analytic functions, we know that there is a neighborhood of the point
  $x_1$ in which both functions coincide with the sum of their Taylor series.
  But having the same derivatives at $x_1$, the two functions also have the same
  Taylor series and thus coincide in a neighborhood of $x_1$. This is absurd
  because, as we defined $x_1$, there must be points arbitrarily close to $x_1$
  where $f(x)\neq g(x)$.
  \end{proof}